\documentclass{patmorin}
\listfiles
\usepackage{pat}
\usepackage{paralist}
\usepackage[T1]{fontenc}
\usepackage[utf8]{inputenc}
% \usepackage{paralist}
\usepackage[inline]{enumitem}
\usepackage{bbm}  % needed for \mathbbm{1}
% \usepackage{logix}
\usepackage{halloweenmath}
\usepackage{stmaryrd}
\usepackage{todonotes}
\usepackage{tcolorbox}
\usepackage{booktabs}
\usepackage{multirow}
\usepackage{comment}

\usepackage{thm-restate}


\newcommand{\dtorso}[2]{{#1}\langle{#2}\rangle^{\downarrow}}
\newcommand{\torso}[2]{{#1}\langle{#2}\rangle}
\DeclareMathOperator{\depth}{depth}


%% so that cleveref outputs Theorems 3--5
\newcommand{\crefrangeconjunction}{--}

% etoolbox allows for robust commands that don't need \protect, e.g.
% \newrobustcmd{\onesub}{\mathord{\includegraphics{figs/one-sub}}}
% \subsection{Approximate Voronoi Diagrams in $G^{\onesub}_k$}
\usepackage{etoolbox}

% david proposes the following additions
\renewcommand{\ge}{\geqslant}
\renewcommand{\le}{\leqslant}
\renewcommand{\geq}{\geqslant}
\renewcommand{\leq}{\leqslant}

\newcommand{\david}[1]{{\color{orange} DW: #1}}
\newcommand{\vida}[1]{{\color{pink} VD: #1}}
\newcommand{\pat}[1]{\textcolor{Blue}{PaMo: #1}}
\newcommand{\gwen}[1]{\textcolor{Purple}{GJ: #1}}
\newcommand{\cyril}[1]{\textcolor{Green}{CG: #1}}
\newcommand{\piotr}[1]{\textcolor{red}{PiMi: #1}}

% \numberwithin{equation}{lem}


\newenvironment{clmproof}{\noindent\emph{Proof of Claim:}}{\hfill\rule{1ex}{1ex}}

\usepackage[longnamesfirst,numbers,sort&compress]{natbib}
\setlength{\bibsep}{0.4ex plus 0.2ex minus 0.2ex}

\usepackage[mathlines]{lineno}
\setlength{\linenumbersep}{2em}
% \linenumbers
% \rightlinenumbers
% \linenumbers
\newcommand*\patchAmsMathEnvironmentForLineno[1]{%
 \expandafter\let\csname old#1\expandafter\endcsname\csname #1\endcsname
 \expandafter\let\csname oldend#1\expandafter\endcsname\csname end#1\endcsname
 \renewenvironment{#1}%
    {\linenomath\csname old#1\endcsname}%
    {\csname oldend#1\endcsname\endlinenomath}}%
\newcommand*\patchBothAmsMathEnvironmentsForLineno[1]{%
 \patchAmsMathEnvironmentForLineno{#1}%
 \patchAmsMathEnvironmentForLineno{#1*}}%
\AtBeginDocument{%
\patchBothAmsMathEnvironmentsForLineno{equation}%
\patchBothAmsMathEnvironmentsForLineno{align}%
\patchBothAmsMathEnvironmentsForLineno{flalign}%
\patchBothAmsMathEnvironmentsForLineno{alignat}%
\patchBothAmsMathEnvironmentsForLineno{gather}%
\patchBothAmsMathEnvironmentsForLineno{multline}%
}



% Taken from
% https://tex.stackexchange.com/questions/42726/align-but-show-one-equation-number-at-the-end
\newcommand\numberthis{\addtocounter{equation}{1}\tag{\theequation}}

\definecolor{brightmaroon}{rgb}{0.76, 0.13, 0.28}
\definecolor{linkblue}{rgb}{0, 0.337, 0.227}
\newcommand{\defin}[1]{\emph{\textcolor{brightmaroon}{#1}}}
\makeatletter
\def\mathcolor#1#{\@mathcolor{#1}}
\def\@mathcolor#1#2#3{%
  \protect\leavevmode
  \begingroup
    \color#1{#2}#3%
  \endgroup
}
\makeatother
\newcommand{\mathdefin}[1]{\mathcolor{brightmaroon}{#1}}
% \newcommand{\mathdefin}[1]{\color{brightmaroon}#1}}
\setlength{\parskip}{1ex}

% Document-specific commands and math operators
\DeclareMathOperator{\tw}{tw}
\DeclareMathOperator{\pw}{pw}
\DeclareMathOperator{\bw}{bw}
\DeclareMathOperator{\td}{td}
%\DeclareMathOperator{\rtw}{rtw}
\DeclareMathOperator{\diam}{diam}
\DeclareMathOperator{\mindist}{min-dist}
\DeclareMathOperator{\dist}{dist}
\DeclareMathOperator{\ld}{ld}
\DeclareMathOperator{\polylog}{polylog}
\DeclareMathOperator{\evol}{Evol}
\DeclareMathOperator{\ivol}{Ivol}
\DeclareMathOperator{\tvol}{Tvol}
\newcommand{\NN}{\mathbb{N}}
\newcommand{\GG}{\mathcal{G}}

\newtheorem{fact}{Fact}

\title{\MakeUppercase{\boldmath Adjacency labelling for minor-closed graph families}}


\author{
 Vida Dujmovi{\'c}\,\footnote{School of Computer Science and Electrical Engineering, University of Ottawa, Ottawa, Canada (\texttt{vida.dujmovic@uottawa.ca}). Research supported by NSERC and a University of Ottawa Research Chair.}
 \qquad
 Gwena\"el Joret\footnote{D\'epartement d'Informatique, Universit\'e libre de Bruxelles, Belgium ({\tt gwenael.joret@ulb.be}). G.\ Joret is supported by the Belgian National Fund for Scientific Research (FNRS) and by the Australian Research Council.}
 \qquad
 Cyril Gavoille\qquad
 Piotr Micek\footnote{Department of Theoretical Computer Science, Jagiellonian University, Kraków, Poland (\texttt{piotr.micek@uj.edu.pl}). Research supported
 the National Science Center of Poland under grant UMO-2018/31/G/ST1/03718 within the BEETHOVEN program.}
 \qquad
 Pat Morin\footnote{School of Computer Science, Carleton University, Ottawa, Canada (\texttt{morin@scs.carleton.ca}). Research supported by NSERC.}
 \qquad
 David~R.~Wood\footnote{School of Mathematics, Monash University, Melbourne, Australia (\texttt{david.wood@monash.edu}). Research supported by the Australian Research Council.}
 }

\date{}


\begin{document}

\maketitle

% \gwen{Another option for the title: Planar graphs as blowups of fans}
%
% \pat{Planar graphs in blowups of fans(?)}
%
% \david{I think `Planar graphs as blowups of fans' reads better than `Planar graphs in blowups of fans', whereas the latter is more accurate. I would weakly lean towards `Planar graphs in blowups of fans'}
%
% \pat{Near rootish blowups of fans contain all planar graphs}
% \david{what does rootish mean, or is it another joke?} \pat{No, not a joke.  Just a way to avoid a square root symbol in a title.} \david{when I read this, I had no idea ``rootish'' refered to square root of $n$.}


\begin{abstract}
  We show that, for every integer $t\ge 1$, there exists a $(1+o(1))\log n$-bit adjacency labelling scheme for the class of $K_t$-minor-free graphs. Equivalently, for every integer $t\ge 1$, there exists an $n^{1+o(1)}$-vertex graph $U_{n,t}$ that contains every $n$-vertex $K_t$-minor-free graph as an induced subgraph.
\end{abstract}

\newpage
\tableofcontents

\newpage
\section{Introduction}

Let $f:\N\to\N$.  A graph family $\mathcal{G}$ \defin{admits an $f(n)$-bit adjacency-labelling scheme} if there is a function $\tau:(\{0,1\}^*)^2\to\{0,1\}$ such that, for each $n\in\N$, each $n$-vertex graph $G$ in $\mathcal{G}$, there exists a function  $\lambda:V(G)\to\{0,1\}^*=\lambda_{G}$ such that,
\begin{itemize}
  \item for each $v,w\in V(G)$, $\tau(\lambda_\omega(v),\lambda(w))=1$ if and only if $vw\in E(G)$; and
  \item for each $v\in V(G)$, $|\lambda(v)|\le f(n)$.
\end{itemize}
The function $\tau$ is called an \defin{adjacency tester} and the function $\lambda$ is called a (vertex) \defin{labelling} of $G$.

A graph $U$ is (induced) \defin{universal} for a set $S$ of graphs if, for each $G\in S$, $U$ contains an induced subgraph that is ismorphic to $G$.  Let $p:\N\to\N$. A graph family $\mathcal{G}$ \defin{admits $p(n)$-vertex universal graphs} if, for each positive integer $n$, there exists a $p(n)$-vertex graph $U_n$ that is universal for the set of $n$-vertex graphs in $\mathcal{G}$.   A classic observation is the following equivalence between adjacency labelling schemes and universal graphs:

\begin{fact}
  A graph family $\mathcal{G}$ admits an $f(n)$-bit adjacency-labelling scheme if and only if $\mathcal{G}$ admits $2^{f(n)}$-vertex universal graphs.
\end{fact}

Several minor-closed classes of graphs admit $(1+o(1))\log n$-bit adjacency labelling schemes, including forests, graphs of bounded treewidth, planar graphs, graphs of bounded genus, and (most generally) apex-minor-free graphs:\footnote{A graph $H$ is an \defin{apex graph} if $H$ contains a vertex $v$ such that $H-v$ is planar.  A graph family $\mathcal{G}$ is \defin{apex-minor-free} if there exists an apex graph $H$ such that no graph in $\mathcal{G}$ contains $H$ as a minor.}

\begin{thm}[\citet{dujmovic.esperet.ea:adjacency}]\label{apex_minor_free}
  For each fixed apex graph $H$, the class of $H$-minor-free graphs
  \begin{itemize}[nosep,nolistsep]
    \item admits a $(1+o(1))\log n$-bit adjacency labelling scheme; and

    \item admits $n^{1+o(1)}$-vertex universal graphs.
  \end{itemize}
\end{thm}

In addition to these, several ``beyond-planar'' graph classes including several non-minor-closed generalizations of $k$-planar graphs admit $(\log n+o(\log n))$-bit adjacency labelling schemes.  Many of these results are based on an adjacency labelling scheme due to \citet{dujmovic.esperet.ea:adjacency} for graphs that have so-called ``product structure.''  Indeed, the labelling schemes for planar, bounded genus, apex-minor-free, and all of the ``beyond planar'' graph classes follow immediately from the fact that all graphs in these graph classes have product structure.

One collection of graph classes that is conspicuously absent from this list is the class of $K_t$-minor-free graphs for any $t\ge 6$.  Indeed, if such a result is true for every fixed positive integer $t$, then this would remove the requirement, in \cref{apex_minor_free}, that the graph $H$ be an apex graph.\footnote{Since $K_t$ contains every $t$-vertex graph as a minor, every $H$-minor-free graph is a $K_{|V(H)|}$-minor-free graph.  Any result for the class of $K_t$-minor-free graphs implies the result for $H$-minor-free graphs for any $t$-vertex graph $H$.}
The best known result in this direction is obtained by combining a $(1+o(1))\log n$-bit labelling scheme of \citet{gavoille.labourel:shorter} for bounded treewidth graphs with a result of \citet{devos.ding.ea:excluding} which states that edges of each $K_t$-minor-free graph can be partitioned into two graphs, each having bounded treewidth:

\begin{thm}[\citet{gavoille.labourel:shorter} with \citet{devos.ding.ea:excluding}]\label{old_best}
  For each fixed graph $H$, the class of $H$-minor-free graphs
  \begin{itemize}[nosep,nolistsep]
    \item admits a $(2+o(1))\log n$-bit adjacency labelling scheme; and

    \item admits $n^{2+o(1)}$-vertex universal graphs.
  \end{itemize}
\end{thm}

\citet{dujmovic.esperet.ea:adjacency} state the problem of reducing the constant from $2$ to $1$ in \cref{old_best} as an open problem.  Unlike apex-minor-free graphs, $H$-minor-free graphs do not necessarily have product structure, which rules out the use of the most powerful tool in the adjacency-labelling toolbox.  In the current work, we resolve this open problem:

\begin{thm}\label{main_result}
  For each fixed graph $H$, the class of $H$-minor-free graphs
  \begin{itemize}[nosep,nolistsep]
    \item admits a $(1+o(1))\log n$-bit adjacency labelling scheme; and

    \item admits $n^{1+o(1)}$-vertex universal graphs.
  \end{itemize}
\end{thm}

\pat{Say something about the proof techniques}

\section{Building Blocks}


Let $T$ be a rooted tree and let $r$ denote the root of $T$. For each $x\in V(T)$, the \defin{$x$-to-root path} \defin{$P_T(x)$} is the path in $T$ from $x$ to $r$.  For each $x\in V(T)$, the \defin{$T$-depth} \defin{$\depth_T(x)$}  is the number of edges in $P_T(x)$.  For each $i\in\N$, the \defin{$i$-th layer} of $T$ is $\mathdefin{L_i(T)}:=\{x\in V(T):\depth_T(x)=i\}$.  For each $x\in V(T)$ and each $a\in V(P_T(x))$ (including $x$),  the vertex $a$ is a $T$-ancestor of $x$ and $x$ is a \defin{$T$-descendant} of $a$.  For each $x\in V(T)$, \defin{$T_x$} denotes the subtree of $T$ induced by all $T$-descendants of $x$ (including $x$).


A \defin{tree decomposition} of a graph $G$ is a collection $(B_x:x\in V(T))$ of vertices subsets of $G$ indexed by the vertices of a tree $T$ such that
\begin{enumerate*}[label=(\alph*)]
  \item for each $v\in V(G)$, $T[\{x\in V(T):v\in B_x\}]$ is a non-empty connected subtree of $T$; and
  \item for each $vw\in E(G)$, $\{v,w\}\subseteq B_x$ for some $x\in V(T)$.
\end{enumerate*}
The \defin{adhesion-width} of a tree decomposition $(B_x:x\in V(T))$ is $\max\{|B_x\cap B_y|:xy\in E(T)\}$.  For a tree decomposition $(B_x:x\in V(T))$ of a graph $G$, the \defin{torso} of $x\in V(T)$, denoted \defin{$\torso{G}{B_x}$} is the supergraph of $G[B_x]$ obtained by adding an edge between $v$ and $w$ for each $\{\{v,w\}\in\binom{B_x\cap B_y}{2}$ and each $y\in N_T(x)$.\footnote{The notation $\torso{G}{B_x}$ introduces some minor ambiguity since the definition of $\torso{G}{B_x}$ depends on $(B_y:y\in N_T(x))$, but the notation makes no reference to the tree decomposition $(B_z:z\in V(T))$.  When we use this notation, the tree decomposition will always be obvious from context.}

When $P$ is path, a tree-decomposition $(B_x:x\in V(P))$ of a graph $G$ is called a \defin{path-decomposition} of $G$.  In this case, we simplify notation by assuming that $P:=1,2,\ldots,r$ and write ``$B_1,\ldots,B_r$ is a path decomposition of $G$.''


\subsection{Degenerate Labelling Schemes}

A graph $G$ is $d$-degenerate if every subgraph of $G$ contains a vertex of degree at most $d$. A graph class $\mathcal{G}$ is $d$-degenerate if every graph in $\mathcal{G}$ is $d$-degenerate. We make use of the following result, which was first used by ?? to find $O(\log n)$-bit adjacency labelling schemes for planar graphs:

\begin{thm}\label{degenerate_labelling_scheme}
  For each $d\in\N$, there exists a $\lceil((d+1)\log n\rceil$-bit adjacency labelling scheme for the class of $d$-degenerate graphs.
\end{thm}




\subsection{Shannon and Shannon-Fano-Elias Codes}

% In most of the following we'll be sloppy and make labels for vertices that have several parts and not discuss how an adjacency-testing procedure can split  a label into its various parts.  This is ok, because in the end the labels we make this way will have parts whose total length is at most $\ell:= \log n+ o(\log n)$ and that are made up of at most $\epsilon\ell$ parts.  This means that we can encode the splitting into parts using $\ell\cdot O(H(\epsilon))$ bits, where $H(\epsilon)=\epsilon \log (1/\epsilon)+(1-\epsilon)\log(1/1-\epsilon)\le 2\epsilon\log(1/\epsilon)$ for $\epsilon\le 1/4$. We will get to choose $\epsilon$ to be decreasing function of $n$.  More on this at the end.

A \defin{code} for a set $S$ is a function $\lambda:S\to\{0,1\}^*$ that maps each element in $S$ to a binary string.  A code is \defin{prefix-free} if $\lambda(x)$ is not a prefix of $\lambda(y)$ for any distinct $x,y\in S$.

We make use of the following result on Shannon Codes:


\begin{lem}\label{shannon_code}
  Let $S$ be a non-empty set, let $\omega:S\to\R^+$, and let $\Omega:=\sum_{i=1}^r\omega(i)$.  Then there exists a prefix-free code $\pi:S\to\{0,1\}^*$ such that $|\pi(x)|\le \log\Omega-\log\omega(i)+1$ for each $x\in S$.
\end{lem}


% \subsection{Vertex-Disjoint Unions of Graphs}
%
%
% The following lemma is not strictly necessary. In fact, it is a special case of \cref{path_family}, below.  We include this lemma and its proof because it provides a good warm-up and highlights some of subtleties of \cref{path_family} .
%
% \begin{lem}\label{vertex_disjoint_union}
%   Let $\mathcal{G}$ be a graph class that admits a $(\log n+g(n))$-bit adjacency labelling scheme.  Let $\mathcal{G}^+$ be the class of graphs that contains a graph $G$ if and only if $G$ is the union of pairwise vertex-disjoint members of $\mathcal{G}$.  Then $\mathcal{G}^+$ admits a $(\log n+1+g(n))$-bit labelling scheme.
% \end{lem}
%
% \begin{proof}
%   Let $G$ be an $n$-vertex graph in $\mathcal{G}^+$ and let $G_1,\ldots,G_r$ be pairwise vertex-disjoint graphs in $\mathcal{G}$ such that $G:=\bigcup_{i=1}^r G_i$.  For each $i\in\{1,\ldots,r\}$, let $\omega(i)=|V(G_i)|$ and observe that $\Omega:=\sum_{i=1}^r \omega(i)=|V(G)|=n$.  Apply \cref{shannon_code} to $\omega$ to obtain a prefix-free code $\pi$.  For each $i\in\{1,\ldots,r\}$ and each vertex $v$ of $G_i$, the label for $v$ has three parts:
%   \begin{enumerate}[nosep,nolistsep,label=(PS\arabic*)]
%     \item The codeword $\pi(v):=\pi(i)$.  This takes $|\pi(i)|\le \log n - \log\omega(i)+1$ bits.
%     \item The adjacency label $\lambda(v)$ for the vertex $v$ in the graph $G_i$.  Since $G_i\in\mathcal{G}$ $|\lambda(v)|\le\log |V(G_i)|+g(|V(G_i)|)=\log\omega(i)+g(\omega(i))\le \log\omega(i)+g(n)$.
%   \end{enumerate}
%   The total length of the label for $v$ is then $\log n + 1 + g(n)$.
%
%   Given the labels for two vertices $v$ and $w$ of $G$, we can then test if $vw\in E(G)$ as follows:
%   \begin{enumerate}[nosep,nolistsep,label=(TS\arabic*)]
%     \item If $\pi(v)\neq\pi(w)$ then $v$ is a vertex of $G_i$ and $w$ is a vertex of $G_j$ for some $i\neq j$. In this case, we immediately conclude that $vw\not\in E(G)$.
%     \item Otherwise, $v$ and $w$ are both vertices of $G_i$. In this case we can use $\lambda(v)$ and $\lambda(w)$ to test if $vw\in E(G_i)$. \qedhere
%   \end{enumerate}
%
%
% \end{proof}

% \subsection{Alphabetic Binary Trees and Shannon-Fano-Elias Codes}

A \defin{rooted ordered binary tree} is a rooted tree with at least two leaves in which each non-leaf node has two children, a \defin{left} child, and a \defin{right} child.
A rooted ordered binary tree $T$ induces an ordering on the the leaves of $T$ in which the leaf descendants  of the left child of $x$ precede the leaf descendants of the right child of $x$ for each non-leaf node $x$ of $T$.
We make use of the following classical result on Alphabetic Binary Trees:


\begin{lem}\label{shannon_fano_elias_tree}\label{shannon_fano_elias}
  Let $r$ be an integer, let $\omega:\{1,\ldots,r\}\to\R^+$, and let $\Omega:=\sum_{i=1}^r\omega(i)$.  Then there exists a rooted ordered binary tree $T_\omega$ whose leaves, in order, are the integer $\{1,\ldots,r\}$ and such that $\depth_{T_\omega}(i)\le \log\Omega - \log\omega(i) + 2$.
\end{lem}

In coding theory, \cref{shannon_fano_elias} is the basis of the Shannon-Fano-Elias coding scheme, as we now explain.  By assigning each edge of $T_\alpha$ a zero or one depending on whether the edge joins a parent to its left or right child, we obtain an encoding of each root-to-leaf path in $T_\alpha$ as a binary string with the following properties:

\begin{cor}\label{shannon_fano_elias_code}\label{shannon_fano_elias_cor}
  Let $r$ be an integer, let $\omega:\{1,\ldots,r\}\to\R^+$, and let $\Omega:=\sum_{i=1}^r\omega(i)$.  Then there exists a function $\rho:\{1,\ldots,r\}\to \{0,1\}^\star$ such that
  \begin{enumerate}[nosep,nolistsep,label=(\alph*)]
    \item $|\rho(i)|\le \log\Omega - \log\omega(i)+2$ for each $i\in\{1,\ldots,r\}$ and
    \item $\rho(i)$ is lexicographically less than $\rho(j)$ for each $1\le i < j\le r$.
  \end{enumerate}
\end{cor}



\subsection{Separators with Low Pathwidth and Low Path-Adhesion-Width}

% A \defin{path decomposition} of a graph $G$ is a sequence $B_1,\ldots,B_r$ of vertex subsets of $G$ such that
% \begin{enumerate*}[label=(\roman*)]
%   \item for each $vw\in E(G)$, $\{v,w\}\subseteq B_i$ for some $i\in\{1,\ldots,r\}$; and
%   \item for each $1\le i < j\le r$, $B_i\cap B_j=\bigcap_{\ell=i}^j B_\ell$.
% \end{enumerate*}
% The \defin{adhesion-width} of a path decomposition $B_1,\ldots,B_r$ is $\max\{|B_i\cap B_{i+1}|:i\in\{1,\ldots,r-1\}\}$.

%
%
%
% For a path decomposition $B_1,\ldots,B_r$ of a graph $G$, we use the conventions $B_0:=B_{r+1}:=\emptyset$.  Then, for each $i\in\{1,\ldots,r\}$, $\mathdefin{G[B_i]^\downarrow}:=G[B_i\setminus B_{i-1}]$ and \defin{$\torso{G}{B_i}$} (the \defin{torso} of $B_i$) denotes the supergraph of $G[B_i]$ obtained adding all edges between all pairs of vertices in $B_{i-1}\cap B_i$ and all pairs of vertices in $B_{i}\cap B_{i+1}$. We use the notation
% \[
%   \mathdefin{\dtorso{G}{B_i}} = \torso{G}{B_i}-(B_{i-1}\cap B_i)
% \]
% to denote the subgraph of the torso $\torso{G}{B_i}$ induced by all vertices that are not contained in $B_{i-1}$.
%
% \begin{lem}\label{path_family}
%    Let $\mathcal{G}$ be a graph family that admits a $(\log n+ g(n))$-bit adjacency labelling scheme.  Let $k>0$ be an integer and let $\mathcal{G}^\star$ be the graph family that contains a graph $G$ if and only if there exists a path decomposition $B_1,\ldots,B_r$ of $G$ of adhesion-width at most $k$ and such that $G[B_i]^{\downarrow}\in\mathcal{G}$, for each $i\in\{1,\ldots,r\}$.  Then $G^\star$ admits a $(\log n + g(n) + O(1+k\log\log n))$-bit adjacency labelling scheme.
% \end{lem}
%
% \pat{The statement of \cref{path_family} might be too coarse. When we apply it, not all the vertices need to pay the extra $O(k\log\log n)$ term. In our application, each $\torso{G}{B_i}$ we be defined by (a small set of) subtrees of the Robertson-Seymour tree. Only the vertices in the root bag of this subtree will need to pay the extra $O(k\log\log n)$.  Vertices in other bags will not have to pay the extra $O(k\log\log n)$ until reaching some lower level of recursion. }
%
% Before proving \cref{path_family}, we make a few remarks.  First, we observe that setting $k=0$ implies a (slightly weaker) version of \cref{vertex_disjoint_union} that replaces ${}+1$ with ${}+O(1)$.  Second, we point out a subtlety in the definition of $G^\star$ that is easily missed on a first reading and is crucial for dealing with the most difficult part of the problem for proper minor-closed graph families (apex vertices coming from distant-ancestor torsos). For each $i\in\{1,\ldots,r\}$, $\dtorso{G}{B_i}$ is ``well-behaved,'' meaning that it has a $(\log n+g(n))$-bit adjacency labelling scheme.  However, for $j\ge 2$, the graph $G[B_j]$ contains the additional vertices in the adhesion $A_j:=B_j\cap B_{j-1}$. There is absolutely no restriction on the adjacency between vertices in $B_j\setminus A_j$ and the vertices in $A_j$: Each vertex $v$ in $A_j$ can be adjacent to any vertex $w$ in $B_{j}\setminus A_j$.
%
% What makes this challenging is that a vertex in $v\in A_j$ can be a vertex of $G[B_{i}]^\downarrow$ for any $i\in\{1,\ldots,j-1\}$.  To deal with this, the same code $\rho(v)=\rho(i)$ that is used to show that $v$ is a vertex of $G[B_{i}]^\downarrow$ must also be used (in conjunction with $\rho(w)=\rho(j)$ to show that $v\in A_j$. To accomplish this label of $w$ contain an additional $O(\log\log n)$ bits for each $v\in A_j$.
%
%
% \begin{proof}[Proof of \cref{path_family}]
%   Let $G$ be an $n$-vertex graph in $\mathcal{G}^\star$ and let $B_1,\ldots,B_r$ be a path decomposition of $G$ that shows $G\in\mathcal{G}^\star$ and that minimizes $r$.
%   Let $A_1:=\emptyset$ and let $A_i:=B_i\setminus B_{i-1}$, for each $i\in\{2,\ldots,r\}$.  For each $i\in\{1,\ldots,r\}$, let $\omega(i):=|B_i\setminus A_i|=|V(\dtorso{G}{B_i})|$.  Note that the minimality of $r$ implies that $\omega(i)=|B_i\setminus A_i|\ge 1$ for each $i\in\{1,\ldots,r\}$.  Observe that the total weight is
%   \[
%     \Omega:=\sum_{i=1}^r\omega(i)=\sum_{i=1}^r|B_i\setminus A_i|=|V(G)|=n \enspace .
%   \]
%   Apply \cref{shannon_fano_elias,shannon_fano_elias_cor} to $\omega$ let $\rho:\{1,\ldots,r\}\to\{0,1\}^*$ be the encoding of the integers $1,\ldots,r$ obtained from the binary tree $T_\omega$.  Since $\Omega=n$ and $\omega(i)\ge 1$,   $\depth_{T_\alpha}(i)\le \log \Omega-\log\omega(i) + 2 \le \log n -\log 1 + 2=\log n + 2$ for each $i\in\{1,\ldots,r\}$.
%
%   For each $v\in V(G)$, let $a(v):=\min\{i\in\{1,\ldots,r\}:v\in B_i\}$ and let $b(v):=\max\{i\in\{1,\ldots,r\}:v\in B_i\}$.  Observe that $B_i\setminus A_i=a^{-1}(i)$, for each $i\in\{1,\ldots,r\}$.  For ech $v\in V(G)$, let $q(v)$ be the lowest common $T_\omega$-ancestor of $a(v)$ and $b(v)$.  We claim that, for each non-leaf node $x$ of $T_\omega$, $|q^{-1}(x)|\le k$.  Indeed, each non-leaf node $x$ of $T_\omega$ has a left child $x_0$ and a right child $x_1$.  The leaves of $T_{x_0}$ are the integers $1,\ldots,i$ and the leaves of $T_{x_1}$ are the integers $i+1,\ldots,r$.  Thus, for any $v\in q^{-1}(x)$, $a(v)\le i$ and $b(v)\ge i+1$.  Therefore $v\in A_{i+1}$. Therefore $|q^{-1}(x)|\le |A_{i+1}|\le k$.
%
%   For each $i\in\{1,\ldots,r\}$ and each $v\in B_i\setminus A_i$, we assign a multi-part label to $v$ as follows:
%   \begin{enumerate}[nosep,label=(P\arabic*)]
%     \item The label $\rho(v):=\rho(i)=\rho(a(v))$ describing the root-to-$i$ path in $T_\omega$.  This label has length at most $\log\Omega-\log\omega(i)+2=\log n-\log\omega(i)+2$.
%     \item The label $\lambda(v)$ assigned to $v$ by the labelling scheme for $\dtorso{G}{B_i}$ which is in the class $\mathcal{G}$. This label has length $\log(|B_i\setminus A_i|)+g(|B_i\setminus A_i|)=\log \omega(i)+g(\omega(i))$.
%     \item A single bit $\nu(v)$ that indicates whether or not $a(v)=b(v)$.
%     \item\label{q_depth} (If $a(v)\neq b(v)$) the integer $\depth_{T_\alpha}(q(v))$.  The maximum depth of any node in $T_\alpha$ is at most $\log n+2$, so this part of the label has length $O(\log\log n)$.
%     \item\label{q_colour} (If $a(v)\neq b(v)$) an integer $\varphi(v)\in\{1,\ldots,k\}$ that is distinct from $\varphi(w)$ for each $w\in q^{-1}(q(v))$.  This part of the label has length $O(\log k)$.
%     \item\label{apex_up} For each $a\in A_i$, the triple
%     \[
%       [\depth_{T_\alpha}(q(a)),\quad \varphi(a),\quad b(v,a)]
%     \]
%     where $b(v,a)$ is a single bit indicating whether or not $va\in E(G)$.  This part of the label has length $O(\log\log n+ \log k + 1)$.  Since $|A_i|\le k$, this part of the label has size $O(k\log\log n)$.
%   \end{enumerate}
%   Thus, in total, the total length of the label for $v$ is $\log n+g(\omega(i)+O(k\log\log n)\le \log n + g(n) + O(k\log\log n)$.
%
%   We now show that these labels are sufficient to perform adjacency testing for $G$.  Given the labels of two vertices $v$ and $w$ in $G$ we proceed as follows:
%   \begin{enumerate}[nosep,label=(T\arabic*)]
%     \item\label{in_torso_test} If $\rho(v)=\rho(w)$ then $v$ and $w$ are both vertices of $\dtorso{G}{B_i}$. In this case, we test if $v$ and $w$ are adjacent using $\lambda(v)$ and $\lambda(w)$.
%     \item\label{apex_test} Otherwise, we may assume without loss of generality that $\rho(v)$ is lexicographically less than $\rho(w)$.  This implies that $v\in B_i\setminus A_i$ and $w\in B_j\setminus A_j$, for some $i < j$.  In this case, the only way in which $v$ and $w$ could be adjacent in $G$ is if $v\in A_{j}$.  Before continuing, we check if $\nu(v)=1$. If not, then $v\not\in A_\ell$ for any $\ell\in\{1,\ldots,r\}$, so $v\not\in A_j$ and we immediately conclude that $vw\not\in E(G)$.  Otherwise $v\in A_\ell$ for at least one $\ell\in\{1,\ldots,r\}$ and the label of $v$ has non-empty parts \ref{q_depth} and \ref{q_colour}.  If $v\in A_{j}$ then \ref{apex_up} of the label for $w$ will contain the triple $(\depth_{T_\alpha}(q(v)),\varphi(v),b(v,w))$. However, the existence of this triple is not sufficient to guarantee that $vw\in E(G)$ since $T$ may have many nodes of depth $d:=\depth_{T_\alpha}(q(v))$ and only one of these is the node $q(v)$. To eliminate this ambiguity, we verify that the length-$d$ prefix of $\rho(v)$ and the length-$d$ prefix of $\rho(w)$ match.
%
%     To summarize, we iterate through the (at most $k$) parts of \ref{apex_up} of the label of $w$ in order to test the following conditions for each $a\in A_{j}$:
%     \begin{enumerate}[nosep,nolistsep]
%       \item Let $d:=\depth_{T_\alpha}(a)$. Is $\depth_{T_\alpha}(q(v))=d$?
%       \item Are first $d$ bits of $\rho(w)$ equal to the first $d$ bits of $\rho(v)$?
%       \item Does $\varphi(a)=\varphi(v)$?
%       \item Is $b(v,a)=1$?
%     \end{enumerate}
%     The answer to all of the first three questions is yes if and only if $a=v$. The answer to the last question is yes if and only $va\in E(G)$.  Thus, the answer to all four questions is yes for exactly one $a\in A_j$ if and only if $vw\in E(G)$.
%   \end{enumerate}
%   Since every edge of $G$ incident to $w$ is an either an edge of $\dtorso{G}{B_{a(w)}}$ (check by \ref{in_torso_test}) or an edge $va$ with $a\in A_i$ (checked by \ref{apex_test}), this correctly determines if $vw\in E(G)$ using only the labels for $v$ and $w$. \end{proof}
%






  \begin{lem}\label{low_pathwidth_separator}
    Let $T$ be a rooted tree, let $\omega:V(T)\to\R^+$ be a weight function and let $\Omega=\sum_{x\in V(T)}\omega(x)$.  For any $0<\alpha< 1$, there exists $X_0\subseteq V(T)$ with the following properties:
    \begin{enumerate}[nosep,nolistsep,label=(\alph*)]
      \item\label{connected} $T[X_0]$ is connected and contains the root of $T$;
      \item\label{skinny} $|X_0\cap L_i(T)|\le 1/\alpha$, for each $i\in\N$; and
      \item\label{light} $\omega(C):=\sum_{y\in V(C)}\omega(y)\le \alpha\Omega$ for each component $C$ of $T-X_0$.
    \end{enumerate}
  \end{lem}

  \begin{proof}
    For each subgraph $C$ of $T$, let $\omega(C):=\sum_{y\in V(C)}\omega(y)$.
    Define $X_0:=\{x\in V(T): \omega(T_x)>\alpha\Omega\}$. We now show that $X_0$ satisfies the conditions of the lemma.

    \noindent\ref{connected}:  Since $\alpha <1$, $X_0$ contains the root $r$ of $T$, since $\omega(T_r)=\omega(T)=\Omega>\alpha\Omega$.  For each $y\in X_0\setminus\{r\}$, $X_0$ contains the $T$-parent $x$ of $y$ since $\omega(T_x)\ge \omega(x)+\omega(T_y)>\omega(T_y)>\alpha\Omega$.  Therefore $T[X]$ contains the root $r$ of $T$ and, for each non-root node $y\in X$, $X$ contains the parent of $y$. Thus $T[X]$ is connected.

    \noindent\ref{skinny}:
    Let $X_i:=X\cap L_i(T)$ for each $i\in\N$.  For each $i\in\N$, the subtrees in $\{T_x:x\in X_i\}$ are pairwise vertex-disjoint, so $\Omega \ge \sum_{x\in X_i} \omega(T_x)> |X_i|\alpha\Omega$.  Rewriting this inequality gives $|X_i|<1/\alpha$ for each $i\in\N$.

    \noindent\ref{light}:
    Each component $C$ of $T-X$ is a subtree of $T$ that is rooted at some node $y$ whose $T$-parent is in $X$.  Since  $y\not\in X$, $\omega(C)=\omega(T_y)\le \alpha\Omega$.
  \end{proof}

  Let $\mathcal{T}:=(B_x:x\in V(T))$ be a rooted tree decomposition of some graph $G$.  A function $\kappa:E(G)\to 2^{V(G)}$ is a \defin{clique assignment function} for $\mathcal{T}$ if, for each edge $xy$ of $T$ where $x$ is the $T$-parent of $y$, $\kappa(xy)$ is the vertex of a maximal clique $K$ in $\torso{G}{B_x}$ whose vertex set contains $B_x\cap B_y$.  When $x\in V(T)$, $\kappa$ is a clique assignment function for $\mathcal{T}$, and $K$ is a clique with $V(K)\subseteq V(G)$, define
  \[
    Y_{\kappa}(x,K) := \{y\in V(T):\text{$x$ is the $T$-parent of $y$ and  $\kappa(xy)=K$}\}
  \]

  \begin{lem}\label{low_adhesion_separator}
    Let $k\in\N$, let $G$ be an $n$-vertex graph containing no clique on $k+1$ vertices, let $\mathcal{T}:=(B_x:x\in V(T))$ be a rooted tree decomposition of $G$ of adhesion-width at most $k$, and let $\kappa$ be a clique assignment function for $\mathcal{T}$.  Then, for any $\alpha\in(0,1)$ there exists a set $X\subseteq V(T)$ such that
    \begin{enumerate}[label=(\alph*)]

      \item\label{x_connected} $T[X]$ is connected and contains the root of $T$;

      \item\label{x_adhesion} $\mathcal{P}:=\left(\bigcup_{x\in L_0(T)\cap X} B_x\right)_{i\in\N}$ is a path decomposition of $G[\bigcup_{x\in X}B_x]$ of adhesion-width less than $2k/\alpha$; and

      \item\label{x_alpha} $\left|\bigcup_{y\in Y_{K,\kappa}(x)\setminus X}\bigcup_{z\in V(T_y)}B_z\setminus V(K)\right|\le \alpha n$ for each $x\in X$ and each maximal clique $K$ in $\torso{G}{B_x}$.
    \end{enumerate}
  \end{lem}

  \begin{proof}
    Let $G$, $\mathcal{T}:=(B_x:x\in V(T))$, and $\kappa$ be defined as in the statement of the lemma.
    Let $r$ be the root of $T$. Define $\omega(r):=|B_r|$ and, for each $y\in V(T)\setminus \{r\}$, define $\omega(y):=|B_y\setminus B_x|$, where $x$ is the with $T$-parent of $y$.  Let $X_0$ be the subset of $V(T)$ obtained by applying \cref{low_pathwidth_separator} to $T$ and $\omega$.
    For each subgraph $C$ of $T$, let $\omega(C)=\sum_{y\in V(C)}\omega(y)$.  Observe that, for each $y\in V(T)$, $\omega(T_y)=|\bigcup_{z\in V(T_y)}B_z\setminus (B_x\cap B_y)|$, where $x$ is the $T$-parent of $y$.  Therefore, $X_0$ contains each $x\in V(T)$ such that $\omega(T_x)>\alpha n$.



    For each $x\in X_0$ and each maximal clique $K$ in $\torso{G}{B_x}$, define \[
       \omega(x,K):=\left|\bigcup_{y\in Y_{\kappa}(x,K)\setminus X_0}\medspace \bigcup_{z\in V(T_y)}B_z\setminus V(K)\right| = \sum_{y\in Y_{\kappa}(x,K)\setminus X_0}\omega(T_y)  \enspace .
    \]
    Let
    \[
      X := X_0\cup\bigcup_{x\in X_0}\{y\in N_T(x)\setminus X_0:\omega(x,K)>\alpha n,\, \kappa(xy)=K\} \enspace .
    \]
    We now show that $X$ satisfies the conditions of the lemma.

    \noindent\ref{x_connected}: $T[X]$ contains the root of $T$ because $X_0$ contains the root of $T$. $T[X]$ is connected because $T[X_0]$ is connected and each vertex $y\in X_0\setminus X$ has a $T$-parent in $X_0$.

    \noindent\ref{x_adhesion}:  For each $i\ge 0$, let $B_i:=\bigcup_{x\in L_{i-1}(T)\cap X} B_x$ and let $r:=\max\{i\ge 1:B_i\neq\emptyset\}$.  We must show that $B_1,\ldots,B_r$ is a path decomposition of $G$ of adhesion-width at most $2k/\alpha$. That $B_1,\ldots,B_r$ is a path decomposition of $G[\bigcup_{x\in X} B_x]$ follows from the fact that $T$ is a tree decomposition of $G$ and that, for each $i\ge 2$ and each $y\in X\cap L_i(T)$, the $T$-parent $x$ of $y$ is in $X\cap L_{i-1}(T)$.

    Now we show that the adhesion-width of $B_1,\ldots,B_r$ is less than $2k/\alpha$.  Let $A_1:=\emptyset$ and let $A_i:=B_{i}\cap B_{i-1}$ for each $i\in\{2,\ldots,r\}$. Fix $i\in\{1,\ldots,r\}$.  We must show that $|A_i|<2k/\alpha$.  Let $v\in A_i$. Then $v$ is contained in $A_y=B_x\cap B_y$, where $x\in X_0\cap L_{i-1}(T)$, $y\in L_i(T)$ and $xy$ is an edge of $T$. We say that $v$ is of Type~0 if $y\in X_{0}$ and $v$ is of Type~1 if $y\in X\setminus X_0$.  The number of Type~0 vertices in $A_i$ is less than $k/\alpha$ since the number of nodes in $X_{0}\cap L_i(T)$ is less than $1/\alpha$ and $|A_y|\le k$ for each $y\in X_{0}$.

    All that remains is to show that the number of Type~1 vertices in $A_i$ is at most $k/\alpha$. Say that a $(x,K)$ is a \defin{heavy pair} if $x\in X\cap L_{i-1}(T)$, $K$ is a maximal clique in $\torso{G}{B_x}$ and $\omega(x,K)>\alpha n$. We will show that the number of heavy pairs is less than $1/\alpha$. This is sufficient, since each Type~1 vertex in $A_i$ belongs to a maximal clique $K$ in $\bigcup_{x\in X\cap L_{i-1}(T)}\torso{G}{B_x}$ with $\omega(x,K)>\alpha n$, and each maximal clique in $\torso{G}{B_x}$ has size at most $k$.

    Let $\mathcal{K}$ be the set of heavy pairs. For each $(x,K)\in\mathcal{K}$, define $C_{x,K}:=\bigcup_{y\in Y_{\kappa}(x,K)}\bigcup_{z\in V(T_y)}B_z\setminus V(K)$.  Then $|C_{x,K}|=\omega(x,K)>\alpha n$.  Furthermore, for two distinct heavy pairs $(x,K),(x',K')\in \mathcal{K}$, $C_{x,K}$ and $C_{x',K'}$ are disjoint.  Therefore,
    \[
      n\ge\left|\bigcup_{(x,K)\in\mathcal{K}} C_{x,K}\right| = \sum_{(x,K)\in\mathcal{K}} \omega(x,K) > |\mathcal{K}|\alpha n \enspace .
    \]
    Dividing both sides by $\alpha n$ shows that $|\mathcal{K}| < 1/\alpha$.

    \noindent\ref{x_alpha}: Assume, for the sake of contradiction, that
    \begin{equation}
       \left|\bigcup_{y\in Y_{K,\kappa}(x)\setminus X}\medspace \bigcup_{z\in V(T_y)}B_z\setminus V(K)\right|> \alpha n \label{contra}
    \end{equation}
    for some $x\in X$ and some maximal clique $K$ in $\torso{G}{B_x}$.  Then, $\omega(T_x) \ge \omega(x,K) > \alpha n$, so $x\in X_0$. Furthermore,
    \[
      \omega(x,K)
       = \left|\bigcup_{y\in Y_{\kappa}(x,K)\setminus X_0}\medspace\bigcup_{z\in V(T_y)}B_z\setminus V(K)\right|
       \ge \left|\bigcup_{y\in Y_{\kappa}(x,K)\setminus X}\medspace\bigcup_{z\in V(T_y)}B_z\setminus V(K)\right|
       > \alpha n \enspace . \\
    \]
    Then, by definition, every vertex in $Y_{\kappa}(x,K)$ is in $X$, so
    \[
      \left|\bigcup_{y\in Y_{K,\kappa}(x)\setminus X}\medspace\bigcup_{z\in V(T_y)}B_z\setminus V(K)\right|=|\emptyset|=0 \enspace ,
    \]
    contradicting \eqref{contra}.
  \end{proof}





  \subsection{The Graph Minor Structure Theorem}

  \pat{Copy definition of $k$-almost-embeddable graph here.  Make sure that the definition stipulates that the maximum clique size is $k$ and that these graphs are $k$-degenerate.}


  \begin{thm}[Graph Minor Structure Theorem]\label{gmst}
    Let $\mathcal{G}$ be a proper minor-closed family of graphs. Then there exists $k:=k_\mathcal{G}$ such that each graph $G\in\mathcal{G}$ has a tree decomposition $\mathcal{T}$ of adhesion-width at most $k$ in which each torso of $\mathcal{T}$ is a $k$-almost embeddable graph.
  \end{thm}



  \subsection{Weighted Adjacency and Clique-Encoding Labelling Schemes}

  Let $g:\N\to\R$.  A graph family $\mathcal{G}$ admits a \defin{weighted $g(n)$-waste adjacency-labelling scheme} if there is a function $\tau:(\{0,1\}^*)^2\to\{0,1\}$ such that, for each $n\in\N$, each $n$-vertex graph $G$ in $\mathcal{G}$, and each assignment $\omega:V(G)\to\R^+$ of positive weights to the vertices of $G$, there exists a function  $\lambda:V(G)\to\{0,1\}^*=\lambda_{G,\omega})$ such that,
  \begin{itemize}
    \item for each $v,w\in V(G)$, $\tau(\lambda_\omega(v),\lambda(w))=1$ if and only if $vw\in E(G)$; and
    \item for each $v\in V(G)$, $|\lambda(v)|\le \log \Omega-\log \omega(v) + g(n)$, where $\Omega=\sum_{x\in V(G)} \omega(x)$ is the total weight of $w$.
  \end{itemize}
  The function $\tau$ is called an \defin{adjacency tester} and the the function $\lambda$ is called a \defin{labelling} of $G$.

  A labelling $\lambda$ of a graph $Q$ is \defin{$p$-clique-encoding} if there exists a function $\tau':(\{0,1\}^*)^3\times\N\to\{0,1\}$ such that, for each maximal clique $K\subseteq Q$, there exists a bitstring $\mu(K)$ of length at most $p$ and an ordering $v_1,\ldots,v_k$ of $V(K)$ such that, for each $w\in V(G)$ and each $i\in\N$, $\tau'(\lambda(v_1),\lambda(w),\mu(K),i)=1$ if and only if $w=v_i$.   In words, the label $\lambda(v_1)$ and the additional bitstring $\mu(K)$ contain enough information to recognize $\lambda(v_i)$ for each $i\in\{1,\ldots,k\}$.

  We will make use of the following result, which can be obtained using the techniques of \citet{dujmovic.esperet.ea:adjacency}.


  \begin{thm}\label{weighted}
    Let $\langle G\rangle$ be a $k$-almost embeddable $n$-vertex graph and let $G$ be a spanning subgraph of $\langle G\rangle$.  Then there exists a $O(n)$-vertex supergraph $Q$ of $\langle G\rangle$ with clique-number at most $2k$ such that the graph $G_Q$ with vertex set $V(G_Q):=V(Q)$ and edge set $E(G_Q):=E(G)\cap E(Q)$ admits a weighted $o(\log n)$-waste adjacency-labelling scheme that is also an $O(k\log\log n)$-clique-encoding labelling of $Q$.
  \end{thm}

  \pat{Note to coauthors: This is the weakest statement I could come up with that would be sufficient for us.  As a result, it may be more complicated than necessary.  In particular, the supergraph $Q$ and the intermediate graph $G_Q$ may be unnecessary. Perhaps we can just find a labelling that is an adacency labelling for $G$ that is also clique encoding for $\langle G\rangle$.}

  \pat{Note to self: Be careful with the order of quantifiers in \cref{weighted}.  We can't have $Q$ depend on the weight function $\omega$ that we choose for $G_Q$.}

  \section{The Proof}

  \begin{proof}[Proof of \cref{main_result}]
    Let $\mathcal{G}$ be a proper minor-closed graph class.  Let $G$ be an $n$-vertex graph in $\mathcal{G}$ and let $\mathcal{T}:=(B_x:x\in V(G))$ be a tree decomposition of $G$ described by \cref{gmst}. Root $T$ at an arbirary vertex $r$.  Define $A_r:=\emptyset$ and, for each edge $xy$ of $T$ where $x$ is the parent of $y$, define $A_y:=B_x\cap B_y$.  Since $\mathcal{T}$ comes from an application of \cref{gmst}, $|A_y|\le k$ for each $y\in V(T)$.

    Let $g(n)\in o(\log n)$ be a function that describes the waste in the labelling scheme of \cref{weighted} and define $\mathdefin{\alpha}:=2^{-\sqrt{(g(n)/k)\log n}}$.
    % The reason for this choice of $\alpha$ is that
    % \[
    %     k\log(1/\alpha)=\frac{g(n)\log n}{\log(1/\alpha)} = \sqrt{k\cdot g(n)\log n}\in o(\log n) ,
    % \]
    % for any fixed $k$.
    What follows below is a proof by induction on $n$.  Within this inductive proof, the value of $\alpha$ does not change, even when the number of vertices in the inductive subproblem decreases.

    We introduce a slightly stronger inductive hypothesis. We show the existence of a fixed constant $c>0$ and an adjacency labelling scheme for $G$ with labelling $\lambda:V(G)\to\{0,1\}^*$ that has the following property: For each $y\in V(T)$ and each $v\in B_y\setminus A_y$,
    \begin{align}
      |\lambda(v)|
      & \le \log n + \frac{g(n)\log n}{\log (1/\alpha)} + c|A_y|\left(\log (k/\alpha) +k\log\log n\right) + ck\log(k/\alpha) \label{solution} \\
      & \le \log n + \frac{g(n)\log n}{\log (1/\alpha)} + ck\left(2\log (k/\alpha) +k\log\log n\right) \label{solution} \notag \\
      & = \log n + O\left(\sqrt{k\cdot g(n)\log n} + k^2\log\log n\right) \notag
      \enspace ,
    \end{align}
    where $c$ is a constant, independent of $n$.
    % This leads to a maximum label length of
    % \begin{align*}
    %    & \log n + O\left(\sqrt{k\cdot g(n)\log n} + k\log k\right) + ck\log\log n \\
    %    & = \log n + o(\log n) + O(k\log k) \enspace .
    % \end{align*}
    Some explation of \cref{solution} is in order:  Let $y$ be a node of $T$ and let $w$ be a vertex in $B_y\setminus A_y$. The label for $w$ has three parts.
    \begin{enumerate}
      \item The first part of $w$'s label is obtained by a recursive procedure that, given a subgraph of size $n'$ recurses on vertex-disjoint subgraphs each of size at most $\alpha n'$.  Thus, $w$ appears in at most $\log_{1/\alpha} n=(\log n)/\log(1/\alpha)$ levels of recursion.  Each level of recursion extends the label of $v$ with an \defin{optimal part} and an \defin{overhead part} of length $g(n)$.  The lengths of the optimal parts of $w$'s label sum to $\log n$ (the first term in \eqref{solution}) and there are $\log_{1/\alpha} n$ overhead parts each of length at most $g(n)$ (the second terms in \eqref{solution}).

      \item The second part of $w$'s label is used by the adjacency testing algorithm to recognize the vertices in $A_y$, each of which may be adjacent to $w$. For each vertex $v\in A_y$, the second part of $w$'s label includes $c(\log(k/\alpha)+k\log\log n)$ bits to help recognize $v$. Thus, the second part of $w$'s label has length at most $c|A_y|(\log(k/\alpha)+k\log\log n)$ (the third term in \eqref{solution}).

      \item The recursive labelling procedure has a base case, which occurs when a subgraph has size at most $k/\alpha$.  This base case is dealt with using \cref{degenerate_labelling_scheme} along with the fact that $G$ is $k$-degenerate. This adds an additional $O(k\log(k/\alpha))$ bits to $w$'s label (the fourth term in \eqref{solution})
    \end{enumerate}

    We first consider the base case of the labelling procedure, which occurs when the input graph has size $n\le k/\alpha$. To handle this case, we make use of the fact that $G$ is $k$-degenerate. By \cref{degenerate_labelling_scheme}, $G$ has a $O(k\log(k/\alpha))$-bit adjacency labelling scheme.

    We now assume that $n> k/\alpha$.
    % For any graph $H$, let $\mathcal{K}(H)$ denote the set of cliques in $H$.\todo{Put this somewhere earlier}.
    We will start by applying \cref{low_adhesion_separator} to the tree decomposition $\mathcal{T}$ of $G$ but this requires a clique assignment function $\kappa$, which we now describe.  For each $x\in X$, we will eventually apply \cref{weighted} to the $k$-almost-embeddable graphs $\torso{G}{B_x}-A_x$ and $G[B_x]-A_x$ to get a labelling $\lambda_x:V(Q_x)\to\{0,1\}^*$ of some supergraph $Q_x$ of $\torso{G}{B_x}-A_x$ that allows for adjacency testing in $G[B_x]-A_x$ and that is $O(k\log\log n)$-clique-encoding for $Q_y$. For each clique $K$ in $\torso{G}{B_x}$, there is a maximal clique $\kappa'(K)$ in $Q_x$ that contains $K$.  For each $T$-child $y$ of $x$, the vertices in $A_y$ form a clique $K_y$ in $\torso{G}{B_x}$. We will use the clique mapping function $\kappa$ defined by $\kappa(xy)=\kappa'(K_y)\cap \torso{G}{B_x}$, for each edge $xy$ of $T$.

    % For each $x\in V(T)$ consider the application of \cref{weighted} to the graph $G[B_x]-A_x$ with supergraph $Q_x:=\torso{G}{B_x}-A_x$.  For each $x\in V(T)$ and each $T$-child $y$ of $x$, $K_y:=Q[A_y\setminus A_x]$ is a clique.  Set $\kappa(y,K_y)$ to be a maximal clique in $Q$ that contains $K_y$. \pat{Come back to this.}

    Apply \cref{low_adhesion_separator} to the graph $G$ with tree decomposition $\mathcal{T}$ with clique assignment function $\kappa$ to obtain $X\subseteq V(T)$ and a path decomposition $B^-_1,\ldots,B^-_r$  of $G[\bigcup_{x\in X} B_x]$ of adhesion-width at most $2k/\alpha$.  For each $x\in X$, let $T^-_x$ be the component of $T-(X\setminus\{x\})$ that contains $x$.  For each $i\in\{1,\ldots,r\}$, define $X_i:=X\cap L_i(T)$ and define $B_i:=\bigcup_{x\in X_i}\bigcup_{y\in T^-_x} (B_y\setminus A_y)$. Let $A_1:=\emptyset$ and, for each $i\in\{2,\ldots,r\}$, let $A_i:=B_i\setminus B_{i-1}$.
    Observe that $B_{i-1}\cap B_i=B^-_{i-1}\cap B^-_{i}$ for each $i\in\{2,\ldots,r\}$, so $B_1,\ldots,B_r$ is a path decomposition of $G$ of adhesion-width at most $2k/\alpha$.

    We now describe the parts of the labels used in the labelling scheme for $G$.  For each $i\in\{1,\ldots,r\}$, let $\omega(i)=|B_i\setminus A_i|$.  Note that $\sum_{i=1}^r\omega(i) = n$.  Apply \cref{shannon_fano_elias_tree,shannon_fano_elias_code} to $\{1,\ldots,r\}$ with weights defined by $\omega$. Let $T_\omega$ be the resulting binary tree and $\rho:\{1,\ldots,r\}\to\{0,1\}^*$ be the resulting prefix-free code.

    \begin{enumerate}[nosep,nolistsep,label=(P\arabic*)]
      \item\label{rho_part} For each $i\in\{1,\ldots,r\}$ and each $v\in B_i\setminus A_i$, the label for $v$ contains a part $\rho(v):=\rho(i)$. This part has length at most $\log n-\log\omega(i) + 2$.
    \end{enumerate}

    Let $v\in V(G)$. Let $a(v):=\min\{i\in\{1,\ldots,r\}:v\in B_i\}$ and let $b(v):=\max\{i\in\{1,\ldots,r\}:v\in B_i\}$.  Then $a(v)$ and $b(v)$ are leaves in $T_\omega$.  Define $q(v)$ to be the lowest common ancestor of $a(v)$ and $b(v)$ in $T_\omega$.  We claim that  $|q^{-1}(x)|< 2k/\alpha$ for each internal node $x$ of $T_\omega$.  Indeed, for each internal node $x$ of $T_\omega$, if $i$ is the minimum leaf in the subtree rooted at the right child of $x$ then the left child of $x$ is the root of a subtree that contains the leaf $i-1$.  Then, every vertex in $q^{-1}(x)$ is in the set $A_i$.  Therefore $|q^{-1}(x)|\le |A_i|< 2k/\alpha$.

    \begin{enumerate}[nosep,nolistsep,label=(P\arabic*)]\setcounter{enumi}{1}
      \item\label{flag_part} For each $v\in\bigcup_{x\in X} B_x$, the label for $v$ contains one bit $t(v)$ that indicates if $a(v)=b(v)$.
      \item\label{lca_part} For each $v\in V(G)$, if $a(v)\neq b(v)$, then the label of $v$ contains $\depth_{T_\omega}(q(v))$. This part has length $O(\log\log n)$.
      \item\label{colour_part} For each $v\in V(G)$, if $a(v)\neq b(v)$, then the label of $v$ contains an integer $\varphi(v)$ that is distinct from $\varphi(w)$ for each $w\in q^{-1}(q(v))$.  This part has length $\lceil\log |q^{-1}(q(v))|\le \lceil \log (2k/\alpha)\rceil$.
    \end{enumerate}

    For each $i\in\{1,\ldots,r\}$ and each $x\in X_i$, define $\omega_i(x):=\omega(T^-_x)$. Note that $\sum_{x\in X_i}\omega_i(x)=|B_i\setminus A_i|=\omega(i)$. For each $i\in\{1,\ldots,r\}$,  apply \cref{shannon_code} with weights defined by $\omega_i$, and let $\pi_i:X^+_i\to\{0,1\}^*$  be the resulting prefix-free code.

    \begin{enumerate}[nosep,nolistsep,label=(P\arabic*)]\setcounter{enumi}{4}
      \item\label{pi_part} For each $i\in\{1,\ldots,r\}$, each $x\in X^+_i$ and each $v\in B_x\setminus A_x$, the label for $v$ contains a part $\pi(v):=\pi_i(x)$. This part has length at most $\log \omega(i)-\log\omega_i(x)$.
    \end{enumerate}

    For each $x\in X$, we now apply \cref{weighted} to the $k$-almost-embeddable graphs $\torso{G}{B_x}-A_x$ and $G[B_x]-A_x$ to get a labelling $\lambda_x:V(Q_x)\to\{0,1\}^*$ of some supergraph $Q_x$ of $\torso{G}{B_x}-A_x$ that allows for adjacency testing in $G[B_x]-A_x$ and that is $O(k\log\log n)$-clique-encoding for $Q_y$. For each clique $K$ in $\torso{G}{B_x}$, there is a maximal clique $\kappa'(K)$ in $Q_x$ that contains $K$.


    % For each $T$-child $y$ of $x$, the vertices in $A_y$ form a clique $K_y$ in $\torso{G}{B_x}$. We will use the clique mapping function $\kappa$ defined by $\kappa(xy)=\kappa'(K_y)\cap \torso{G}{B_x}$, for each edge $xy$ of $T$.
    %
    %
    %
    % \Cref{weighted} provides a supergraph $Q_x$ of $\torso{G}{B_x}$ and a mapping $\zeta$ that maps each clique $K$ in $\torso{G}{B_x}$ onto a clique $\zeta(K)$ in $Q_x$ that contains $K$.


    We now explain how to assign weights to vertices of $Q_x$ to use in \cref{weighted}. For each $x\in X$ and each clique $K$ in $\torso{G}{B_x}-A_x$, let $Y^-_{\kappa}(x,K):=Y_{\kappa}(x,K)\setminus X$ and let $\omega^{-}(x,K):=\sum_{y\in Y'_{\kappa}(x,K)}\omega(T_{y})$.  For $x\in X$, and each $v\in V(Q_x)$, let $\mathcal{K}_v:=\{K\in\mathcal{K}(\torso{G}{B_x}-A_x):v\in V(\kappa'(K))\}$. Finally, for each $x\in X$ and each $v\in V(Q_x)$, define $\omega_x(v):=1+\sum_{K\in\mathcal{K}_v}\omega(x,K)$.
    For each $x\in X$, each clique in $Q_x$ contains at most $2k$ vertices, so $\sum_{v\in B_x\setminus A_x}\omega_x(v)=|V(Q_x)|+\sum_{v\in B_x\setminus A_x}\sum_{K\in\mathcal{K}_v}\omega^-(x,K)\le O(|B_x\setminus A_x|)+ k\sum_{K\in\mathcal{K}_x}\omega^-(x,K)\le (2k+O(1))\omega_i(x)$.

    % For each $x\in X$, let $Q_x:=\torso{G}{B_x}-A_x$.
    Now we apply the weighted labelling scheme in \cref{weighted} to the graphs $\torso{G}{B_x}$ and $G[B_x]-A_x$ with vertex weights defined by $\omega_x:V(Q_x)\to\N$.  This gives a weighted adjacency labelling $\lambda_x$ of the intermediate graph $G_{Q_x}$. The labelling $\lambda_x$ is also $O(k\log\log n)$-clique encoding $Q_x$.  Since the intermediate graph $G_{Q_x}$ contains $G[B_x]-A_x$ as an induced subgraph, $\lambda_x$ is also an adjacency labelling of $G[B_x]-A_x$.
    % We now distinguish between two kinds of vertices in $G$.
     % contains the following parts:

    \begin{enumerate}[nosep,nolistsep,label=(P\arabic*)]\setcounter{enumi}{5}
      \item\label{dejmmw_part} For each $i\in\{1,\ldots,r\}$, each $x\in X_i$ and each $v\in B_x\setminus A_x$, the label for $v$ contains $\lambda_x(v)$.

      For each $i\in\{1,\ldots,r\}$, each $x\in X$, each $y\in N_T(x)\setminus X$, each $z\in V(T_y)$, and each $w\in B_z-A_z$, the label for $w$ contains $\lambda(\kappa'(\kappa(xy)))$.  More precisely, the label for $w$ contains $\lambda_x(v_1)$ and $\mu(\kappa'(\kappa(xy)))$ as in the definition of clique-encoding labellings.  This part has length at most $O(\log (k\cdot\omega_i(x)-\log\omega(v))+g(n)$.

      \item\label{recursive_part} For each $x\in X$, each maximal clique $K$ in $\torso{G}{B_x}$ the subgraph $G_{x,K}:=G[\bigcup_{y\in Y_{x,K}\setminus X}\medspace\bigcup_{z\in V(T_y)} (B_z\setminus A_z)]$ is labelled recursively.  For each vertex $v\in V(G_{x,K})$ this part of the label is called $R(v)$.  By the inductive hypothesis,
      \begin{align*}
        |R(v)|
        & \le \log |V(G_{x,K})|+\frac{g(n)\log(\alpha n)}{\log(1/\alpha)}+ c|A_z\setminus A_x|\log\log n \\
        & = \log\omega(x,K) + \frac{g(n)\log n}{\log(1/\alpha)}-g(n)+ c|A_z\setminus A_x|\log\log n \\
        & \le \log\omega(w) + \frac{g(n)\log n}{\log(1/\alpha)}-g(n)+ c|A_z\setminus A_x|\log\log n \enspace ,
      \end{align*}
      where $w$ is any vertex in $V(K)$.
    \end{enumerate}
    Finally, after doing any recursive labelling in \ref{recursive_part}, we augment some of the labels.


    \begin{enumerate}[nosep,nolistsep,label=(P\arabic*)]\setcounter{enumi}{7}
      \item\label{apex_edges} For each $x\in X$ and, each $z\in V(T^-_x)$ and each $w\in B_z\setminus A_z$, the label for $w$ will contain a part for each vertex $v\in A_z$.

      \begin{enumerate}
        \item If $v\in A_x$, then $a(v)<i$ and $b(v)\ge i$. Therefore $q(a(v))$ is a $T_\omega$-ancestor of $i$.  In this case, the label for $w$ contains $\depth_{T_\alpha}(q(a(v))$ and $\varphi(v)$.

        \item If $v\in B_x\setminus A_x$ and $z= x$ then $v$ and $w$ are both vertices in $G[B_x]-A_x$. In this case, \ref{dejmmw_part} already contains enough information to test the presence of the edge $vw$.

        \item If $v\in B_x\setminus A_x$ and $z\neq x$, then $z$ is a strict $T$-descendant of $x$.  Let $y$ be the $T$-child of $x$ that is a $T$-ancestor of $z$.  In this case, part \ref{dejmmw_part} of the label for $w$ contains the label for the clique $K_y$ that contains $v$.  We augment the label for $w$ with $O(\log\log n)$ bits indicating this, and $O(\log k)$ bits needed to identify $v$ from \ref{dejmmw_part}.

        \item If $v\not\in B_x$, then $v$ and $w$ are both vertices of $G_{x,K}$ for some clique $K$ in $\torso{G}{B_x}-A_x$.  In this case, \ref{recursive_part} contains enough information to the presence of the edge $vw$.
      \end{enumerate}
    \end{enumerate}

    Before discussing the adjacency testing procedure, we first analyze the lengths of the labels. For each $x\in X$ and each $v\in V(T^-_x)$ the total length of parts \ref{flag_part}--\ref{dejmmw_part} is
    \[
      \log n - \log\omega(i) + \log\omega(i) - \log\omega_i(x) + \log(k\cdot\omega_i(x))-\log\omega(x,K) + \log\omega(x,K) + ....
    \]
  \end{proof}



    % In this labelling, there are $\log_{1/\alpha} n$ levels of recursion, where level $i$ consists of subgraphs each of size at most $\alpha^i n$.  At each layer, there is a waste of $O(\sqrt{\log n})$.  When a vertex stops, there is a one-time waste (from \ref{colour_part}) of $\log(2k/\alpha)$. So the label length is
    % \[
    %   \log n + \frac{\log n\cdot \sqrt{\log n}}{\log 1/\alpha} + \log(1/\alpha)
    % \]
    % If we take $\alpha=2^{-(\log n)^{3/4}}$ then this becomes
    % \[
    % \log n + \frac{\log n\cdot \sqrt{\log n}}{(\log n)^{3/4}} + (\log n)^{3/4} = \log n + O((\log n)^{3/4})
    % \]


  \bibliographystyle{plainurlnat}
  \bibliography{mflabelling}

\end{document}
